\section{String: Literal, Escape, dan Metode Dasar}

String adalah urutan karakter yang merepresentasikan teks. Di JavaScript, string dapat ditulis dengan kutip tunggal (\code{'}), kutip ganda (\code{"}), atau backtick (\code{`}) untuk template literal. Template literal (ES6) mendukung interpolasi (\code{`Hello \$\{name\}`}) dan multi-baris. String bersifat \textbf{immutable}---metode string mengembalikan string baru tanpa mengubah aslinya \cite{jsInfo}.

\begin{table}[htbp]
\centering
\begin{tabular}{|P{3.5cm}|P{4cm}|P{4.5cm}|}
\hline
\textbf{Metode} & \textbf{Fungsi} & \textbf{Contoh} \\
\hline
\code{.length} & Jumlah karakter & \code{"Hello".length} $\rightarrow$ 5 \\
\hline
\code{.toUpperCase()} & Huruf besar & \code{"hi".toUpperCase()} $\rightarrow$ \code{"HI"} \\
\hline
\code{.includes(str)} & Cek substring & \code{"Hello".includes("ell")} $\rightarrow$ true \\
\hline
\code{.indexOf(str)} & Posisi substring & \code{"Hello".indexOf("l")} $\rightarrow$ 2 \\
\hline
\code{.slice(start, end)} & Ambil bagian & \code{"Hello".slice(1,4)} $\rightarrow$ \code{"ell"} \\
\hline
\code{.split(sep)} & Pecah ke array & \code{"a,b".split(",")} $\rightarrow$ \code{["a","b"]} \\
\hline
\code{.trim()} & Hapus spasi ujung & \code{" Hi ".trim()} $\rightarrow$ \code{"Hi"} \\
\hline
\code{.replace(old,new)} & Ganti substring & \code{"hi".replace("h","H")} $\rightarrow$ \code{"Hi"} \\
\hline
\end{tabular}
\caption{Metode string JavaScript yang sering digunakan}
\end{table}

Escape sequence digunakan untuk karakter khusus dalam string: \code{\\n} (baris baru), \code{\\t} (tab), \code{\\\\} (backslash), \code{\\"} (kutip ganda). Dengan template literal, escape sequence jarang dibutuhkan karena multi-baris dan interpolasi sudah didukung secara native \cite{mdnJsGuide}.

\begin{catatan}
\textbf{Template Literal untuk Kode Bersih.} Alih-alih konkatenasi (\code{"Hello " + name + "!"}), gunakan template literal (\code{`Hello \$\{name\}!`}). Selain lebih mudah dibaca, template literal mendukung ekspresi (\code{`Total: \$\{price * qty\}`}) dan multi-baris tanpa \code{\\n}. Ini adalah salah satu fitur ES6 yang paling bermanfaat dalam praktik sehari-hari \cite{jsInfo}.
\end{catatan}

